% Section 2 -- Related work (condensed 2026-11-23). Each subsection: what the line of work does, what it assumes, where this work differs.
\section{Related work}\label{sec:related}

\subsection{Site effects and harmonization}

Harmonization removes the shift and scaling that each site adds to the features. ComBat~\cite{johnson2007,fortin2018} and its deep variants~\cite{an2025deepresbat} protect biological variables such as the diagnosis by including them as covariates, which works only if the protected variable varies within sites. When groups are unevenly distributed across batches, batch correction can shrink true group differences or create spurious ones~\cite{nygaard2016}, ComBat that protects the diagnosis leaks label information when class proportions differ across sites~\cite{nieto2026}, residual site effects can bias accuracy unless the evaluation accounts for sites~\cite{solanes2021}, and when the covariate distributions of two batches barely overlap, causal methods decline to decide whether a batch effect exists~\cite{bridgeford2025}. Traveling-subject designs separate measurement from recruitment differences~\cite{yamashita2019} but are costly and rare. Class-pure sites are the case of no overlap at all. Our design step creates the missing overlap, and our estimator does not need to estimate site effects.

\subsection{Shortcut learning and debiasing}

Deep models often adopt shortcuts that work on standard benchmarks but fail to transfer~\cite{geirhos2020}; in medical imaging the shortcut is often the hospital or acquisition source~\cite{zech2018,degrave2021}. Remedies include group reweighting~\cite{sagawa2020}, contrastive learning across attributes~\cite{zhang2022cnc}, margin-based emphasis of bias-conflicting samples~\cite{wang2025mdn} and class-wise loss balancing~\cite{zuo2026}. All of them remain well defined on class-pure data, but the data lack a contrast between classes within a site, so the diagnostic effect cannot be separated from the site effects without further assumptions (Proposition~\ref{prop:pure}); contrastive and margin-based methods also need bias-conflicting samples, of which there are none. Domain generalization for classes seen in a single domain~\cite{kaai2025fond} transfers invariance learned from classes that appear in several domains; with class-pure sites, invariance learned across sites of the same class leaves the difference between case sites and control sites (Section~\ref{sec:prop3}). Bridging subjects supply the missing bias-conflicting samples, and our within-site loss uses only pairs that share the spurious attribute and differ in label, the hard negatives of contrastive debiasing~\cite{zhang2022cnc}.

\subsection{Multi-site classification, confounded designs and evaluation}

Methods for multi-site classification from functional connectivity were mostly developed on data in which most sites contain both classes, for example site-aware population graphs~\cite{parisot2018,wu2027croip} and domain adaptation~\cite{song2024braindas,fang2025scda}; federated learning~\cite{wang2025gafd}, minimization of site dependence~\cite{kunda2023} and adversarial or mixed-effects removal of site information~\cite{ganin2016,dinsdale2021,zhao2020cfnet,nguyen2023armed} were also developed and evaluated on such data, as was CORE~\cite{wang2026core}, which corrects confounders site by site. We test two of them in Section~\ref{sec:res-pureonly}. In genomics, confounding between batch and group inflates cross-validated performance~\cite{soneson2014}, and when each of two batches held one outcome, a predictor trained on them failed on the unseen combinations of batch and outcome~\cite{parker2012}. Design-based remedies place shared material in several batches: reference panels~\cite{song2020}, remeasured controls~\cite{ye2023} and technical bridging samples~\cite{bridge2022}. Our bridging subjects are instead new individuals of the other class at the same site, with no repeated measurement, so whether they work depends on whether the site effect is additive. For evaluation, leaving out whole sites is an established protocol in multi-site autism classification~\cite{abraham2017}. Our leave-site-pair-out evaluation applies leave-pair-out cross-validation, which gives nearly unbiased AUC estimates in small samples~\cite{airola2011}, to whole sites; because discrimination within clusters differs from discrimination between clusters~\cite{vanklaveren2014}, Section~\ref{sec:eval} states which of the two it estimates. Evaluation protocols can also hide leakage through model selection~\cite{yang2027} or instability across training domains~\cite{zhang2025pf2hc}.
